\scanpage{243}
\mspage{124}

Il serait facile d'expliciter les conditions sur les données (65) (66) pour qu'elles correspondent à un multigraphe à homothéties et gr. d'automs homogène $H$ --- je préfère plutôt \unsure{d\textquoteright{}écrire|écrire} les conditions pour que les données

\[
H,Z\subset\unsure{\mathfrak S|\mathcal G}_S\quad\text{et}\quad (s_0,t_0)\in S\times S
\tag*{(68)}
\]

correspondent à un multigraphe à homothéties et groupe d'autom \deleted{\unsure{pour}}\inserted{hom} --- en notant que la donnée de $(s_0,t_0)\in A$ \deleted{permet}, en plus de $H,Z$, permet de récupérer $A$, comme

\[
A=\bigl\{\{u\unsure{\lambda|(}s_0,u\lambda t_0\}\mid\lambda\in Z,\ u\in\deleted{\unsure{G|O}}H\bigr\}.
\tag*{(69)}
\]

On trouve le système de conditions

\[
\left\{\begin{aligned}
&\text{a) }H\text{ transitif sur }S,\ H\text{ et }Z\text{ commutent.}\\
&\text{b) }\forall\lambda\in Z,\ \exists u\in H\quad\text{tel que}\\
&\qquad u(s_0,t_0)=(\lambda^{-1}s_0,\lambda t_0).
\end{aligned}\right.
\tag*{(70)}
\]

On récupère alors \deleted{$A_0$}

\[
A_0=\bigl\{\{u\deleted{\unsure{s|\lambda}}s_0,ut_0\}\mid u\in H\bigr\},
\tag*{(\unsure{78|71})}
\]

et par b) $Z$ sera un groupe d'homothéties pour $(S,A_0)$, ce qui engendre un multigraphe à homothéties $(S,A)$ :

On n'exclut pas le cas où on aurait $s_0=t_0$, auquel cas $A_0$ est formé de toutes les parties ponctuelles de $S$, et la condition b) signifie que $\lambda\in Z\Longrightarrow\lambda^2=1$. Si $\lambda\in Z$, $\lambda\ne1$, alors $A_\lambda$ est formé des orbites sous $\lambda$, qui sont des parties à deux éléments. \deleted{L'étoile}

% archive sheet number: 242
